\section{Conclusion}
\label{sec:conclusion}

We presented \textbf{Belief-Constrained Inference} (\bci{}), an inference-time
framework that externalizes, verifies, and corrects intermediate visual beliefs in
VLMs before final answer generation. Our core finding is that premise errors
constitute 33.3\% of all incorrect VLM answers on GQA, rising to 50\% for
logical/relational questions, and that this rate is a conservative lower bound due to
verifier coverage gaps in attribute and action claim families.

Correcting these premises yields $+18.0$~pp accuracy gain at $n{=}500$ under a
full-replacement oracle (McNemar $p{=}1.5{\times}10^{-14}$), scaling to $+16.2$~pp
at $n{=}2000$ ($p{=}1.78{\times}10^{-48}$). A practical confidence-gated policy
(E12b) achieves $+3.45$~pp at $n{=}2000$ ($p{=}1.3{\times}10^{-5}$) with
59.85\% intervention rate and multi-seed stability of $66.18\%\pm 0.14$~pp.
Targeted correction outperforms random belief deletion by $+35$~pp, confirming
premise-specificity. Against matched lightweight baselines, \bci-E1 leads the
competitor matrix by $+16.6$~pp, while Self-Correct falls below Vanilla by $10.0$~pp.

\paragraph{Limitations.}
Current results are limited to GQA with deterministic scene-graph verifiers.
The E1 policy is an oracle upper bound requiring ground-truth facts at test time;
E12b is practical but still relies on scene-graph verification.
Verifier coverage is incomplete (49\% mean), especially for action and OCR claims.
Cross-benchmark generalization (MMMU, MathVista, MM-Vet) and model-agnosticity
(InstructBLIP, LLaVA-Next) remain for future work.
Multi-seed stability for E1 at $n{=}2000$ shows zero variance across seeds,
suggesting the sampling procedure may benefit from additional randomization checks.

\paragraph{Future Work.}
The most impactful extensions are:
(1)~replacing the scene-graph verifier with a learned or retrieval-augmented visual
oracle to remove the ground-truth dependency;
(2)~extending evaluation to open-ended generation benchmarks;
(3)~incorporating uncertainty estimation into belief revision;
(4)~conducting a latency/accuracy Pareto analysis to characterize the full operating
frontier of confidence-gated \bci{};
and (5)~cross-benchmark generalization studies to validate premise-error patterns
beyond GQA.

\paragraph{Broader Impact.}
\bci{} provides an interpretable audit trail---the explicit belief set and
verification labels---that makes VLM reasoning more transparent and correctable.
This is particularly valuable in safety-sensitive applications where understanding
\emph{why} a model failed is as important as preventing failure.
